\documentclass[10pt,oneside]{book}
\usepackage[margin=1in]{geometry}
\usepackage{amsmath,amssymb,amsthm,mathtools,mathrsfs}
\usepackage{fontspec}
\setmonofont[Path=fonts/,Scale=0.85]{JetBrainsMono.ttf}
\usepackage{newunicodechar}
\usepackage{booktabs,longtable,array,enumitem}
\usepackage{xcolor}
\usepackage[colorlinks=true,linkcolor=blue!50!black,urlcolor=blue!50!black,citecolor=blue!50!black]{hyperref}
\setlist{itemsep=2pt,topsep=4pt}
\sloppy\emergencystretch=5em

\newtheorem{theorem}{Theorem}[chapter]
\newtheorem{proposition}[theorem]{Proposition}
\newtheorem{lemma}[theorem]{Lemma}
\newtheorem{corollary}[theorem]{Corollary}
\theoremstyle{definition}
\newtheorem{definition}[theorem]{Definition}
\newtheorem{example}[theorem]{Example}
\theoremstyle{remark}
\newtheorem{remark}[theorem]{Remark}

\newcommand{\Lean}[1]{\texttt{#1}}
\newcommand{\leanref}[3]{\par{\footnotesize\raggedright\noindent\textsf{Lean:}\ \texttt{#1}\ \ (\texttt{#2}, lines #3)\par}}
\newcommand{\Q}{\mathbf{Q}}\newcommand{\Z}{\mathbf{Z}}\newcommand{\C}{\mathbf{C}}
\newcommand{\R}{\mathbf{R}}\newcommand{\N}{\mathbf{N}}\newcommand{\F}{\mathbf{F}}
\newcommand{\Sha}{\mbox{\fontencoding{U}\fontfamily{wncy}\selectfont Sh}}
\renewcommand{\Sha}{\mathrm{Sha}}
\renewcommand{\Re}{\operatorname{Re}}\renewcommand{\Im}{\operatorname{Im}}
\DeclareMathOperator{\Gal}{Gal}\DeclareMathOperator{\GL}{GL}\DeclareMathOperator{\U}{U}
\DeclareMathOperator{\Frob}{Frob}
\newcommand{\ASK}{\mathrm{ASK}}\newcommand{\Main}{\mathrm{Main}}
\newcommand{\Ga}{\mathrm{Ga}}\newcommand{\Fo}{\mathrm{Fo}}\newcommand{\Ee}{\mathrm{Ee}}
\newcommand{\sg}{\mathrm{sg}}\newcommand{\gainD}{\mathrm{gainD}}

% common Unicode characters that may appear in running text
\newunicodechar{ψ}{\ensuremath{\psi}}\newunicodechar{χ}{\ensuremath{\chi}}\newunicodechar{ε}{\ensuremath{\varepsilon}}
\newunicodechar{ℓ}{\ensuremath{\ell}}\newunicodechar{≤}{\ensuremath{\le}}\newunicodechar{≥}{\ensuremath{\ge}}
\newunicodechar{→}{\ensuremath{\to}}\newunicodechar{∈}{\ensuremath{\in}}\newunicodechar{×}{\ensuremath{\times}}
\newunicodechar{θ}{\ensuremath{\theta}}\newunicodechar{ρ}{\ensuremath{\rho}}\newunicodechar{σ}{\ensuremath{\sigma}}
\newunicodechar{μ}{\ensuremath{\mu}}\newunicodechar{λ}{\ensuremath{\lambda}}\newunicodechar{π}{\ensuremath{\pi}}
\newunicodechar{Φ}{\ensuremath{\Phi}}\newunicodechar{Λ}{\ensuremath{\Lambda}}\newunicodechar{ζ}{\ensuremath{\zeta}}
\newunicodechar{γ}{\ensuremath{\gamma}}\newunicodechar{δ}{\ensuremath{\delta}}\newunicodechar{η}{\ensuremath{\eta}}
\newunicodechar{α}{\ensuremath{\alpha}}\newunicodechar{β}{\ensuremath{\beta}}\newunicodechar{τ}{\ensuremath{\tau}}
\newunicodechar{ω}{\ensuremath{\omega}}\newunicodechar{ξ}{\ensuremath{\xi}}\newunicodechar{κ}{\ensuremath{\kappa}}
\newunicodechar{≠}{\ensuremath{\ne}}\newunicodechar{∞}{\ensuremath{\infty}}\newunicodechar{⊂}{\ensuremath{\subset}}
\newunicodechar{⊆}{\ensuremath{\subseteq}}\newunicodechar{∑}{\ensuremath{\sum}}\newunicodechar{√}{\ensuremath{\surd}}\newunicodechar{ℝ}{\ensuremath{\mathbb R}}

% Printed contents list parts and chapters; PDF bookmarks keep sections.
\setcounter{tocdepth}{0}
\AtBeginDocument{\hypersetup{bookmarksdepth=2}}
